\documentclass[11pt]{article}
\usepackage[T1]{fontenc}
\usepackage{lmodern}
\usepackage{amsmath,amssymb}
\usepackage[margin=1in]{geometry}
\usepackage[colorlinks=true,linkcolor=blue,citecolor=blue,urlcolor=blue]{hyperref}

\title{Literature Status: Twisted Inclusion and Alternating 2-Cocycles}
\author{Run \texttt{fa\_banach\_001}, agent \texttt{agent\_super\_01}}
\date{2026-07-01}

\begin{document}
\maketitle

\section*{Status}

This is a lightweight \texttt{literature\_already\_answered} packet. It records
that the future-work application announced in Choi,
\emph{A twisted inclusion between tensor products of operator spaces},
arXiv:1606.06287v2, was written up in the separate later paper Choi,
\emph{Constructing alternating 2-cocycles on Fourier algebras},
arXiv:2008.02226v4.

\section*{Source Future-Work Item}

The source note arXiv:1606.06287 proves that for operator spaces $V$ and $W$,
if $\widetilde W$ denotes the opposite operator-space structure on the
underlying Banach space of $W$, then the identity map on $V\otimes W$ extends
to a contraction
\[
  V\widehat{\otimes} W \longrightarrow V\otimes_{\min}\widetilde W .
\]
This is Theorem \texttt{t:mainthm} in the parsed source.

The abstract says that this theorem will be applied in future work to construct
antisymmetric $2$-cocycles on certain Fourier algebras. The author comments
added on 20 May 2020 say that the cocycle application was removed from the
historical note and would be written separately in fuller detail.

\section*{Supporting Answer}

The supporting paper arXiv:2008.02226 is exactly such a separate write-up. Its
abstract states that it gives the first examples of locally compact groups
whose Fourier algebras support non-zero alternating $2$-cocycles, and it lists
the twisted inclusion theorem for operator-space tensor products as one of the
two main technical ingredients. The arXiv metadata for arXiv:2008.02226v4 also
states that the paper includes an improved version of material from
arXiv:1606.06287v2.

The theorem-level evidence is explicit. Theorem \texttt{t:headline} proves
that for $n\geq 4$ and $G$ equal to $\mathrm{SU}(n)$, $\mathrm{SL}(n,\mathbb
R)$, or $\mathrm{Isom}(\mathbb R^n)$, there is a non-zero continuous
alternating $2$-cocycle on the Fourier algebra $A(G)$. Theorem
\texttt{t:twisted inclusion} restates/proves the twisted inclusion. Theorem
\texttt{t:mainthm} constructs non-zero alternating $2$-cocycles from non-zero
co-completely bounded derivations, and Theorem \texttt{t:underyournoses}
supplies the co-cb derivations needed for the headline examples.

\section*{Scope}

This is not a new proof. It records the already-known later answer to the
future-work application announced in arXiv:1606.06287: applying the twisted
inclusion theorem to construct antisymmetric, equivalently alternating,
$2$-cocycles on certain Fourier algebras. New follow-up questions posed in the
final section of arXiv:2008.02226, such as the small-dimensional cases, are not
settled by this status packet.

\section*{References}

\begin{itemize}
\item Y. Choi, \emph{A twisted inclusion between tensor products of operator
spaces}, arXiv:1606.06287v2.
\item Y. Choi, \emph{Constructing alternating 2-cocycles on Fourier algebras},
arXiv:2008.02226v4.
\end{itemize}

\end{document}
